\documentclass[12pt,a4paper]{article}
\usepackage[margin=2.5cm]{geometry}
\usepackage{amsmath}
\usepackage{hyperref}
\usepackage{parskip}

\title{\textbf{Modelling}\\[0.3em]
\large Part V(d) --- Modelling Report\\
Asymmetric Sentiment Effects in Earnings Calls}
\author{NLP Course --- Group Project}
\date{}

\begin{document}
\maketitle

\section{Task Definitions}

We frame the prediction problem in two complementary ways:

\begin{itemize}
  \item \textbf{Regression}: predict CAR$_{[0,3]}$ (a continuous return) from
        transcript features and controls.
  \item \textbf{Classification}: predict whether CAR$_{[0,3]} > 0$
        (binary direction label).
\end{itemize}

The regression task is the primary one because it supports the formal Wald test
of the asymmetry hypothesis.

\section{Regression Models}

\subsection{OLS with HC3 Standard Errors}

The main specification is:
\begin{equation}
  \text{CAR}_{i,[0,3]} = \beta_0 + \beta_1 \, \text{NegRate}_i +
  \beta_2 \, \text{PosRate}_i + \mathbf{x}_i^\top \boldsymbol{\gamma} + \varepsilon_i
\end{equation}
where controls $\mathbf{x}_i$ include pre-event volatility and year dummies.
OLS is estimated with HC3 heteroskedasticity-consistent standard errors via
\texttt{statsmodels}.  The asymmetry test is a Wald test of
$H_0: \beta_1 + \beta_2 = 0$.

\subsection{Ridge Regression}

Ridge ($\ell_2$ penalty, $\alpha = 1$) is used when high-dimensional features
(TF-IDF, year dummies) are included to control overfitting.  With only $\sim$155
training events, plain OLS on 500+ features is severely underdetermined.

\subsection{Random Forest and MLP}

Random Forest (100 trees) and a two-hidden-layer MLP are included as nonlinear
baselines.  Both use \texttt{random\_state=42} for reproducibility.

\subsection{Null Model}

The null model predicts the training-set mean for every test event.
Positive OOS $R^2$ means a model beats this naive baseline.

\section{Classification Models}

Binary label: $y = 1$ if CAR$_{[0,3]} > 0$.
We evaluate five classifiers:

\begin{itemize}
  \item Naive Bayes (Gaussian, for LM features; Multinomial, for TF-IDF)
  \item Logistic Regression (\texttt{max\_iter=1000})
  \item Decision Tree
  \item Random Forest (100 trees)
  \item Multi-Layer Perceptron (two hidden layers)
\end{itemize}

All use \texttt{random\_state=42}.  Class imbalance in the test window
(positive CARs dominated 2019--2020) means F1 and AUC are more informative
than accuracy alone.

\section{Time-Series Cross-Validation}

Beyond the single train/test split, we run an expanding-window cross-validation:
train on 2016$\rightarrow$year $Y$, test on year $Y{+}1$, for $Y \in \{2017, 2018, 2019\}$.
Pooled OOS $R^2$ averages across all folds gives a more robust estimate
than a single test window.

\section{Code Reference}

\begin{itemize}
  \item \texttt{src/nasdaq\_nlp/models/regression.py} --- OLS, Ridge, RF, MLP, Wald test.
  \item \texttt{src/nasdaq\_nlp/models/classifiers.py} --- classification models.
  \item \texttt{src/nasdaq\_nlp/evaluation/} --- OOS $R^2$, MAE, AUC utilities.
  \item \texttt{notebooks/03\_modeling.ipynb} --- full model training and result export.
\end{itemize}

\end{document}
